\documentclass[11pt,a4paper]{article}
\usepackage[utf8]{inputenc}
\usepackage[T1]{fontenc}
\usepackage{amsmath,amssymb,amsthm}
\usepackage{mathtools}
\usepackage{geometry}
\geometry{margin=2.5cm}
\usepackage{hyperref}
\usepackage{booktabs}
\usepackage{enumitem}
\usepackage{amsfonts}
\usepackage{amstext}

\newtheorem{theorem}{Theorem}[section]
\newtheorem{proposition}[theorem]{Proposition}
\newtheorem{lemma}[theorem]{Lemma}
\newtheorem{corollary}[theorem]{Corollary}
\newtheorem{definition}[theorem]{Definition}
\newtheorem{remark}[theorem]{Remark}
\newtheorem{hypothesis}[theorem]{Hypothesis}

\newcommand{\Ceff}{\mathcal{C}_{\mathrm{eff}}}
\newcommand{\Cadm}{\mathcal{C}_{\mathrm{adm}}}

\title{Topological Invariants of Admissible Configurations:\\
Charge Quantisation, Absence of Magnetic Monopoles,\\
and Projective Chirality in Cosmochrony}
\author{J\'er\^ome Beau\\
\small Independent Researcher, France; jerome.beau@cosmochrony.org}
\date{2026}

\begin{document}
  \maketitle

  \begin{abstract}
    We classify the topological invariants of admissible effective configurations
    in the Cosmochrony framework.
    The projection fiber $\Pi \cong S^3$ carries a Hopf fibration
    $S^1 \hookrightarrow S^3 \to S^2$, whose $U(1)$ fiber sector supports
    electric charge as a winding number associated with
    $\pi_1(S^1) \cong \mathbb{Z}$.
    On a canonical admissible configuration model, the bounded-flux constraint
    removes polar regions from the base~$S^2$, reducing the admissible
    base to a topological annulus with $\pi_2 = 0$.
    This triviality excludes isolated magnetic monopoles as admissible
    topological excitations.
    We further argue that parity violation arises from projective chirality:
    a structural orientation dependence of admissibility under non-injective
    projection, which selects one handedness of winding invariants over its
    conjugate.
    We distinguish rigorously between results established on the canonical
    model and open problems related to the full effective configuration
    space~$\Ceff$.
  \end{abstract}

  \smallskip\noindent
  \textbf{Keywords.}
  Topological invariants, admissible configurations, charge quantisation,
  magnetic monopoles, projective chirality, parity violation,
  Hopf fibration, Cosmochrony.

  \clearpage
  \tableofcontents

  \clearpage
  \section{Introduction}
    \label{sec:intro}

    \subsection{Context}

      The Cosmochrony framework~\cite{Cosmochrony} derives physical observables
      from the constrained relaxation dynamics of a pre-geometric relational
      substrate~$\chi$.
      Previous companion papers established:
      (i)~the emergence of electric charge as a structural invariant of
      non-injective projection~\cite{Cosmochrony,Beau2026e};
      (ii)~the dynamical mechanism linking mass and charge to a shared
      relaxation budget under bounded-flux constraints~\cite{Beau2026e};
      (iii)~the role of non-injective projection in generating $U(1)$ gauge
      structure via the Hopf fibration of the projection
      fiber~$\Pi \cong S^3$~\cite{Cosmochrony}.

      These works address the \emph{mechanism} by which charge and mass emerge.
      The present paper addresses a different question:
      \emph{what are the admissible topological classes of effective
      configurations, and what physical constraints do they impose?}

    \subsection{Contribution}

      We identify a topological layer of the Cosmochrony framework,
      complementary to the dynamical and spectral analyses developed
      elsewhere.
      The central results are:
      \begin{enumerate}[label=(\arabic*)]
\item Electric charge corresponds to a $\pi_1$ invariant of the $U(1)$
fiber, with $\pi_1(S^1) \cong \mathbb{Z}$
(Section~\ref{sec:charge}).
\item On the canonical admissible model, the second homotopy group of
the admissible base is trivial, $\pi_2(\Cadm) = 0$, excluding
isolated magnetic monopoles
(Section~\ref{sec:monopoles}).
\item Parity violation can be interpreted as a manifestation of
projective chirality: an orientation dependence of admissibility
under the non-injective projection~$\Pi$
(Section~\ref{sec:chirality}).
      \end{enumerate}

      These three results are unified by a single structural principle:
      the physics of the effective description is governed by the type of
      topological invariants that the projection renders admissible, and
      by the orientation it selects among them.

    \subsection{Scope and limitations}

      Results~(1) and~(2) are established on a canonical admissible
      configuration model (Section~\ref{sec:canonical}).
      Their extension to the full effective configuration space~$\Ceff$
      is an open problem.
      Result~(3) is established at the framework level as a structural
      mechanism; its identification with the specific $V{-}A$ structure of the
      weak interaction remains qualitative.

  \section{Admissible Configuration Space}
    \label{sec:admissible}

    \subsection{The projection fiber}

      The relational substrate~$\chi$ is accessed through admissible local
      projections onto a reduced projection fiber~\cite{Cosmochrony}
      \begin{equation}
        \label{eq:fiber}
        \Pi \cong S^3.
      \end{equation}
      This identification follows from the minimal compact geometry required
      to support non-injective yet stable projections
      (\S5.1 of~\cite{Cosmochrony}).
      The associated Hopf fibration
      \begin{equation}
        \label{eq:hopf}
        S^1 \hookrightarrow S^3 \xrightarrow{\;\pi_H\;} S^2
      \end{equation}
      naturally decomposes~$\Pi$ into a circular fiber~$S^1$ and a
      two-dimensional base~$S^2$.

    \subsection{Physical identification of the fibration sectors}

      The Hopf fibration induces a decomposition of the gauge structure
      of the projected description~\cite{Cosmochrony}:
      the $U(1)$ phase degree of freedom along the $S^1$ fiber
      corresponds to electromagnetism,
      while the $SU(2)$ rotational degrees of freedom of the $S^3$ fiber
      itself correspond to the weak interaction.
      This decomposition is not postulated but emerges from the geometry of
      the projection.
      This identification follows from the correspondence established
      in~\cite{Cosmochrony} (\S5.2) and is not re-derived in the present
      paper; we use it as structural input.

    \subsection{Bounded-flux constraint and admissibility}

      The bounded-flux constraint~\cite{Cosmochrony,Beau2026c}
      \begin{equation}
        \label{eq:flux}
        |\partial_t \chi_v| \leq c_\chi
        \qquad \text{for every node } v \in \mathcal{G}(n),
      \end{equation}
      imposes a maximal relaxation rate on every relational variable.
      At the level of the projected description, this translates into a
      constraint on the admissible gradient of the fiber
      phase~\cite{Beau2026e}:
      \begin{equation}
        \label{eq:fluxproj}
        |\nabla \theta| \leq \frac{c_\chi}{\hbar_\chi},
      \end{equation}
      where $\theta \in [0, 2\pi)$ is the $U(1)$ fiber coordinate
      and $\hbar_\chi$ is the effective action scale.

      \begin{definition}[Admissible configuration space]
        \label{def:Cadm}
        The admissible configuration space~$\Cadm$ is the subspace of
        fiber configurations satisfying the bounded-flux
        constraint~\eqref{eq:fluxproj} at every point.
      \end{definition}

      The constraint~\eqref{eq:fluxproj} excludes configurations with
      arbitrarily large phase gradients.
      Its topological consequence is the subject of the next sections.

  \section{Fundamental Group and Electric Charge}
    \label{sec:charge}

    \subsection{Winding number as topological invariant}

      The $U(1)$ fiber of the Hopf fibration has the topology of a circle.
      Its fundamental group is
      \begin{equation}
        \label{eq:pi1}
        \pi_1(S^1) \cong \mathbb{Z}.
      \end{equation}
      A closed loop in the fiber defines a winding number
      $w \in \mathbb{Z}$, which is a topological invariant of the
      configuration.

    \subsection{Electric charge as winding}

      In the Cosmochrony framework, the bounded relaxation flux
      $\vec{J}_\chi \sim \nabla \chi$ decomposes in the weak-gradient regime
      as~\cite{Cosmochrony}
      \begin{equation}
        \label{eq:decomp}
        \nabla \chi = -\nabla \varphi + A_T,
        \qquad \nabla \cdot A_T = 0,
      \end{equation}
      where $\varphi$ is a scalar potential and $A_T$ is the transverse
      (solenoidal) component.
      The effective charge invariant is~\cite{Cosmochrony}
      \begin{equation}
        \label{eq:charge}
        q_{\mathrm{eff}} = \kappa \oint_\gamma A_T \cdot \mathrm{d}\ell,
      \end{equation}
      which is the circulation of the transverse flux around a closed loop
      encircling the excitation.
      This circulation is an integer multiple of the fundamental winding
      quantum, giving
      \begin{equation}
        \label{eq:quantised}
        q_{\mathrm{eff}} = w \cdot e,
        \qquad w \in \mathbb{Z},
      \end{equation}
      where $e$ is the elementary charge.

      \begin{proposition}[Charge quantisation]
        \label{prop:charge}
        In the Cosmochrony framework, electric charge is quantised as a
        consequence of $\pi_1(S^1) \cong \mathbb{Z}$.
        The winding number~$w$ of a fiber configuration is a topological
        invariant that cannot change under continuous deformation.
        The admissibility constraint~\eqref{eq:fluxproj} removes
        configurations exceeding the flux bound but does not introduce
        identifications that would contract the fiber.
        Therefore the fundamental group of the fiber sector
        remains that of $\mathrm{U}(1)$.
      \end{proposition}

      \begin{remark}[Relation to the Dirac quantisation condition]
        \label{rem:dirac}
        The standard Dirac argument derives charge quantisation from the
        \emph{existence} of magnetic monopoles, via the single-valuedness
        of the wavefunction on a sphere surrounding the monopole.
        In Cosmochrony, charge quantisation has a different origin:
        it follows from the topology of the $U(1)$ fiber itself,
        independently of whether magnetic monopoles exist.
        The two mechanisms are logically independent.
      \end{remark}

    \subsection{Sign of charge and conjugation}

      The winding number~$w$ carries a sign.
      Configurations with $w > 0$ and $w < 0$ correspond to opposite
      chiralities of the transverse torsion pattern.
      Charge conjugation corresponds to the reversal $w \to -w$,
      i.e.\ to a relational conjugation of the fiber
      configuration~\cite{Cosmochrony}.

  \section{Canonical Admissible Model and Triviality of $\pi_2$}
    \label{sec:canonical}

    \subsection{The full Hopf base: $S^2$ and magnetic monopoles}

      In the \emph{unconstrained} Hopf fibration, the base manifold is the
      full two-sphere~$S^2$.
      The second homotopy group of $S^2$ is non-trivial:
      \begin{equation}
        \label{eq:pi2S2}
        \pi_2(S^2) \cong \mathbb{Z}.
      \end{equation}
      A non-trivial element of $\pi_2(S^2)$ would correspond, in the
      gauge-theoretic language, to a Dirac magnetic
      monopole~\cite{Dirac1931}:
      a point-like source of magnetic flux classified by an integer
      topological charge.
      In non-abelian gauge theories, the 't~Hooft--Polyakov
      construction~\cite{tHooft1974,Polyakov1974} generalises this to
      solitonic monopole solutions classified by $\pi_2(G/H)$.

      If the full $S^2$ were the admissible base, magnetic monopoles would
      be topologically admissible.

    \subsection{Bounded-flux reduction of the base}

      The bounded-flux constraint~\eqref{eq:fluxproj} restricts the set of
      admissible fiber configurations.
      At the level of the Hopf base~$S^2$, the constraint excludes
      configurations in which the fiber phase varies too rapidly.

      Consider the canonical admissible model in which the constraint
      removes a neighbourhood of each pole of~$S^2$
      (the north and south poles being the points where the fiber
      degenerates, corresponding to maximal flux concentration).
      The admissible base is then
      \begin{equation}
        \label{eq:annulus}
        S^2_{\mathrm{adm}}
        = S^2 \setminus (D^2_N \cup D^2_S)
        \simeq S^1 \times (0,1),
      \end{equation}
      where $D^2_N$ and $D^2_S$ are open discs around the north and south
      poles, respectively.

      \begin{proposition}[Homotopy of the canonical admissible base]
        \label{prop:homotopy}
        The canonical admissible base~$S^2_{\mathrm{adm}}$ is homotopy
        equivalent to $S^1 \times (0,1)$.
        Its homotopy groups are
        \begin{equation}
          \label{eq:homotopyadm}
          \pi_1(S^2_{\mathrm{adm}}) \cong \mathbb{Z},
          \qquad
          \pi_2(S^2_{\mathrm{adm}}) = 0.
        \end{equation}
      \end{proposition}

      \begin{proof}
        Removing two disjoint open discs from $S^2$ yields a surface
        homeomorphic to a closed annulus, which deformation-retracts onto
        its core circle~$S^1$.
        A circle has $\pi_1 = \mathbb{Z}$ and $\pi_k = 0$ for all $k \geq 2$.
      \end{proof}

    \subsection{Status of the canonical model}

      \begin{remark}[Scope of Proposition~\ref{prop:homotopy}]
        \label{rem:scope}
        The result is established for the canonical model in which the
        admissibility constraint removes exactly two polar discs.
        Its extension to the full admissible configuration space~$\Ceff$
        requires showing that the bounded-flux constraint generically reduces
        the base to a space with trivial~$\pi_2$.
        This is a well-posed topological problem but remains open.
        The expectation that the result extends is motivated by the fact
        that the bounded-flux constraint removes precisely the gradient
        singularities required to support non-trivial $\pi_2$ classes:
        any admissible base must lack the polar regions where the fiber
        degenerates, and it is this degeneration that generates
        $\pi_2(S^2) \cong \mathbb{Z}$.
      \end{remark}

  \section{Absence of Magnetic Monopoles}
    \label{sec:monopoles}

    \subsection{Topological exclusion}

      \begin{corollary}[No isolated magnetic charge]
        \label{cor:nomonopole}
        On the canonical admissible model, no isolated magnetic monopole is
        a topologically admissible excitation.
      \end{corollary}

      \begin{proof}
        An isolated magnetic monopole corresponds to a non-trivial class in
        $\pi_2$ of the gauge bundle base.
        In the present framework, the configuration space already
        incorporates the gauge structure through the projection fiber,
        so that $\pi_2$ of $\Ceff$ directly classifies admissible magnetic
        defects.
        By Proposition~\ref{prop:homotopy},
        $\pi_2(S^2_{\mathrm{adm}}) = 0$,
        so no such class exists.
      \end{proof}

    \subsection{Physical interpretation}

      The absence of magnetic monopoles is one of the most robust empirical
      facts in physics: despite extensive searches across energy scales
      and astrophysical environments, no isolated magnetic charge has been
      observed~\cite{PDG,Patrizii2015}.

      In the standard Dirac--'t~Hooft--Polyakov
      framework~\cite{Dirac1931,tHooft1974,Polyakov1974},
      monopoles are topologically admissible but their absence must be
      explained dynamically (e.g.\ by cosmological dilution in inflationary
      models).

      In Cosmochrony, the situation is reversed:
      the absence of monopoles is not a dynamical accident but a
      \emph{topological consequence of admissibility}.
      The bounded-flux constraint removes the topological support for
      monopole-type excitations.

      \begin{remark}[Precise status]
        \label{rem:monopolestatus}
        In Cosmochrony, the absence of magnetic monopoles is not postulated.
        It is already realised on the canonical admissible configuration model
        by the triviality of~$\pi_2$.
        The remaining open problem is to lift this result from the canonical
        admissible base to the full effective configuration space~$\Ceff$.
      \end{remark}

    \subsection{Compatibility with Dirac strings}

      The triviality of $\pi_2$ on the admissible base does not affect the
      Dirac string construction for \emph{non}-monopolar configurations.
      Magnetic flux lines (corresponding to $B = \nabla \times A_T$) remain
      admissible as long as they form closed loops.
      Only \emph{isolated} magnetic sources, which would require non-trivial
      $\pi_2$ classes, are excluded.

  \section{Projective Chirality and Parity Violation}
    \label{sec:chirality}

    \subsection{Orientation of the winding invariant}

      The winding number $w \in \mathbb{Z}$ carries a sign determined by the
      orientation of the fiber phase relative to a reference direction.
      In a parity-symmetric theory, configurations with $w$ and $-w$ would
      be related by a mirror transformation and would have equal dynamical
      weight.

    \subsection{Projective chirality}

      In Cosmochrony, the projection $\Pi: \chi \to \chi_{\mathrm{eff}}$
      is non-injective: multiple relational configurations project onto the
      same effective description~\cite{Cosmochrony}.
      This non-injectivity can be \emph{chiral}: the projection may
      treat configurations with opposite orientations asymmetrically.

      \begin{definition}[Projective chirality]
        \label{def:chirality}
        A projection is said to be projectively chiral if the stabilisation
        rates satisfy
        \begin{equation}
          \label{eq:chiral}
          \Gamma(Q) \neq \Gamma(-Q),
        \end{equation}
        where $\Gamma(Q)$ is the effective stabilisation rate for configurations
        of topological charge~$Q$, and $-Q$ denotes the charge-conjugate
        configuration.
      \end{definition}

      When the projection is chiral, parity is violated:
      mirror-related configurations are not equally admissible.

    \subsection{Structural mechanisms for chirality selection}

      Two mechanisms within the Cosmochrony framework produce structural
      chirality.

      \paragraph{Spinorial selection.}
        The spectral admissibility programme~\cite{Beau2026a1,Beau2026a2}
        establishes that the bounded-flux constraint selects spinorial
        (spin-$\frac{1}{2}$) representations as the most admissible along the
        chain $Q_8 \subset 2I \subset \mathrm{SU}(2)$.
        Spinorial representations carry an intrinsic handedness through the
        double-cover structure $\mathrm{SU}(2) \to \mathrm{SO}(3)$:
        the central $\mathbb{Z}_2$ retained by binary covers is precisely
        the element that distinguishes left- and right-handed spinors.

      \paragraph{Non-injective projection.}
        The non-injectivity of the projection creates a fiber degeneracy:
        multiple $\chi$-configurations map to the same effective state.
        The orientation of the winding within this fiber is not preserved
        by a generic projection, so parity invariance is not automatic.
        A chiral projection selects one orientation over its conjugate,
        producing $\Gamma(Q) \neq \Gamma(-Q)$ without any explicit
        parity-breaking term in the underlying admissibility structure.

  \subsection{Relation to the observed $V{-}A$ structure}

    The maximal parity violation observed in weak interactions,
    most directly in the polarised beta decay of
    ${}^{60}\mathrm{Co}$~\cite{Wu1957},
    can be interpreted within this framework as the manifestation of a
    maximally chiral projection in the $\mathrm{SU}(2)$ sector of~$\Pi$.

    The spectral admissibility hierarchy provides the structural
    mechanism: the bounded-flux filter selects specific spinorial sectors,
    and the non-injectivity of the projection fixes the orientation.

    \begin{remark}[Status of the $V{-}A$ identification]
      \label{rem:VA}
      This correspondence is currently qualitative.
      A derivation of the \emph{maximal} chirality ($V{-}A$ rather than a
      partial chiral bias) requires showing that the admissibility filter
      selects a single chirality sector completely.
      This is structurally plausible given the strict spinorial selection
      established in the spectral programme, but the explicit identification
      of weak-interaction sectors within Weil blocks remains open.
    \end{remark}

  \section{Unified Topological Picture}
  \label{sec:unified}

  The three results of this paper are unified by a single principle:
  the physics of the effective description is determined by the
  topological invariants that the admissibility constraint renders
  available, and by the orientation that the projection selects.

  \begin{center}
    \begin{tabular}{lll}
      \toprule
      Topological invariant & Physical consequence & Mechanism \\
      \midrule
      $\pi_1(S^1) \cong \mathbb{Z}$ &
      Charge quantisation &
      Winding of $U(1)$ fiber \\
      $\pi_2(S^2_{\mathrm{adm}}) = 0$ &
      No magnetic monopoles &
      Bounded-flux reduction of base \\
      Orientation of $\pi_1$ &
      Parity violation &
      Projective chirality \\
      \bottomrule
    \end{tabular}
  \end{center}

  This table makes explicit that charge, the absence of monopoles,
  and parity violation are not three independent phenomena but three
  aspects of the same topological structure:
  the admissible homotopy groups of the projected fiber and the
  orientation that the projection selects among their generators.

  Spectral admissibility (the O-series~\cite{Beau2026a7,Beau2026a11})
  does not change the topology of~$\Ceff$;
  it selects which topological sectors are dynamically realised
  and determines their mass hierarchy.
  The topological and spectral layers are complementary:
  the first classifies \emph{what} invariants exist,
  the second determines \emph{how} they organise into a spectrum.

  \section{What is Proved and What is Open}
  \label{sec:status}

  \begin{center}
    \begin{tabular}{lll}
      \toprule
      Result & Status & Reference \\
      \midrule
      $\pi_1(S^1) \cong \mathbb{Z}$ (charge quantisation) &
      Proved (canonical model) &
      Prop.~\ref{prop:charge} \\
      $\pi_2(S^2_{\mathrm{adm}}) = 0$ (no monopoles) &
      Proved (canonical model) &
      Prop.~\ref{prop:homotopy} \\
      Projective chirality ($\Gamma(Q) \neq \Gamma(-Q)$) &
      Established (framework) &
      Def.~\ref{def:chirality} \\
      Spinorial selection &
      Proved (spectral programme) &~\cite{Beau2026a1} \\
      \midrule
      $\pi_2(\Ceff) = 0$ in general &
      Open &
      Remark~\ref{rem:scope} \\
      Maximal chirality ($V{-}A$) &
      Open &
      Remark~\ref{rem:VA} \\
      Weak sectors in Weil blocks &
      Open &~\cite{Beau2026a15} \\
      \bottomrule
    \end{tabular}
  \end{center}

  \section{Conclusion}
  \label{sec:conclusion}

  We have identified the topological invariants of admissible effective
  configurations in the Cosmochrony framework and derived their physical
  consequences.

  Electric charge arises from $\pi_1$ invariants of the $U(1)$ Hopf
  fiber, with quantisation following from $\pi_1(S^1) \cong \mathbb{Z}$.
  The bounded-flux constraint reduces the admissible Hopf base from
  $S^2$ to a topological annulus with $\pi_2 = 0$, excluding isolated
  magnetic monopoles as admissible topological excitations.
  Projective chirality provides a structural interpretation of parity
  violation, arising from the orientation dependence of admissibility
  under non-injective projection.

  These three results are unified by a single structural principle:
  the effective physics is governed by the admissible homotopy
  classes of the projected fiber and the orientation selected by the
  projection.

  The open problems are well-posed:
  (i)~extending the triviality of $\pi_2$ from the canonical model to
  the full $\Ceff$;
  (ii)~deriving the maximal ($V{-}A$) chirality from the spectral
  admissibility hierarchy;
  (iii)~identifying the weak-interaction sectors within the Weil-block
  decomposition of the spectral programme.

  This paper establishes a topological layer of the Cosmochrony
  framework, complementary to the dynamical analysis
  of~\cite{Beau2026e} and the spectral analysis of the O-series.
  The three layers together provide a classification (topological),
  a mechanism (dynamical), and a quantification (spectral) of
  the physical observables emerging from constrained
  non-injective projection.

  \begin{thebibliography}{99}

    % White paper
\bibitem{Cosmochrony}
J.~Beau,
``Cosmochrony: A Pre-Geometric Relational Framework for Emergent Physics,''
white paper, Zenodo (2025).
DOI: \href{https://doi.org/10.5281/zenodo.17957509}{10.5281/zenodo.17957509}.

% Born-Infeld
\bibitem{Beau2026c}
J.~Beau,
``Emergent Spacetime Geometry from Bounded Relational Relaxation,''
companion paper, Zenodo (2026).
DOI: \href{https://doi.org/10.5281/zenodo.18407505}{10.5281/zenodo.18407505}.

% SM-Charge-Mass
\bibitem{Beau2026e}
J.~Beau,
``Charge, Mass, and Inertia as Saturated Responses in a Pre-Geometric Relational Description,''
companion paper, Zenodo (2026).
DOI: \href{https://doi.org/10.5281/zenodo.18488987}{10.5281/zenodo.18488987}

% Admissibility
\bibitem{Beau2026a1}
J.~Beau,
``Spectral Admissibility under Bounded Flux: Non-Abelian Mode Selection on Cayley Graphs of SU(2) Subgroups,''
spectral admissibility subprogram paper, Zenodo (2026).
DOI: \href{https://doi.org/10.5281/zenodo.18944985}{10.5281/zenodo.18944985}.

% Capacity
\bibitem{Beau2026a2}
J.~Beau,
``Spectral Capacity Functional and the Binary-Polyhedral Maximality Conjecture,''
spectral admissibility subprogram paper, Zenodo (2026).
DOI: \href{https://doi.org/10.5281/zenodo.18945273}{10.5281/zenodo.18945273}.

% Stratigraphy
\bibitem{Beau2026a4}
J.~Beau,
``Three Particle Generations and Mass Hierarchy from Spectral Stratigraphy,''
spectral admissibility subprogram paper, Zenodo (2026).
DOI: \href{https://doi.org/10.5281/zenodo.19040490}{10.5281/zenodo.19040490}.

% relaxation
\bibitem{Beau2026a5}
J.~Beau,
``Asymptotic Saturation of Projective Resolution: Expander Relaxation Graphs,''
spectral admissibility subprogram paper, Zenodo (2026).
DOI: \href{https://doi.org/10.5281/zenodo.19056975}{10.5281/zenodo.19056975}.

% O1
\bibitem{Beau2026a6}
J.~Beau,
``Projective Resolution Dynamics on Variable-Valence Relaxation Graphs,''
O1 spectral admissibility subprogram paper, Zenodo (2026).
DOI: \href{https://doi.org/10.5281/zenodo.19076319}{10.5281/zenodo.19076319}.

% O3
\bibitem{Beau2026a7}
J.~Beau,
``Projective Dynamics and Mass Hierarchy: Hierarchical Amplification via Growing Relational Valence,''
O3 spectral admissibility subprogram paper, Zenodo (2026).
DOI: \href{https://doi.org/10.5281/zenodo.19076244}{10.5281/zenodo.19076244}.

% O4
\bibitem{Beau2026a8}
J.~Beau,
``Subdiffusive Valence Growth under Bounded Relational Flux: Structural Derivation of the Cascade Exponent $\beta$,''
O4 spectral admissibility subprogram paper, Zenodo (2026).
DOI: \href{https://doi.org/10.5281/zenodo.19101472}{10.5281/zenodo.19101472}.

% O5
\bibitem{Beau2026a9}
J.~Beau,
``Admissible Frontier Saturation and the Cascade Exponent: A Representation-Theoretic Obstruction and its Matrix-Level
Refinement,''
O5 spectral admissibility subprogram paper, Zenodo (2026).
DOI: \href{https://doi.org/10.5281/zenodo.19120328}{0.5281/zenodo.19120328}.

% O6
\bibitem{Beau2026a10}
J.~Beau,
``Matrix-Level Dynamic Redundancy and the Structural Confinement of the Cascade Exponent~$\beta$,''
O6 spectral admissibility subprogram paper, Zenodo (2026).
DOI: \href{https://doi.org/10.5281/zenodo.19140624}{10.5281/zenodo.19140624}.

% O7
\bibitem{Beau2026a11}
J.~Beau,
``Projective Capacity and the Continuum Limit of Admissible Redundancy,''
O7 spectral admissibility subprogram paper, Zenodo (2026).
DOI: \href{https://doi.org/10.5281/zenodo.19149773}{10.5281/zenodo.19149773}.

% O8
\bibitem{Beau2026a12}
J.~Beau,
``Three-Step Path Fingerprints on LPS Graphs and the Capacity Growth Exponent,''
08 spectral admissibility subprogram paper, Zenodo (2025).
DOI: \href{https://doi.org/10.5281/zenodo.19151769}{10.5281/zenodo.19151769}.

% O9
\bibitem{Beau2026a13}
J.~Beau,
``Projective Capacity Beyond Expanders: Polynomial-Growth Relaxation Graphs and the Capacity Exponent,''
09 spectral admissibility subprogram paper, Zenodo (2026).
DOI: \href{https://doi.org/10.5281/zenodo.19154517}{10.5281/zenodo.19154517}

% O10
\bibitem{Beau2026a14}
J.~Beau,
``Polynomial-Growth Capacity Dynamics at Large $q$: Volumetric Growth Confirmation, Capacity Decay, and the Weil
Decomposition Roadmap,''
010 spectral admissibility subprogram paper, Zenodo (2026).
DOI: \href{https://doi.org/10.5281/zenodo.19161074}{10.5281/zenodo.19161074}

% O11
\bibitem{Beau2026a15}
J.~Beau,
``Weil-Block Projective Capacity on Heisenberg Graphs: Resolving the Representation Obstruction and Extracting the
Capacity Exponent,''
011 spectral admissibility subprogram paper, Zenodo (2026).
DOI: \href{https://doi.org/10.5281/zenodo.19169987}{10.5281/zenodo.19169987}

% 012
\bibitem{Beau2026a16}
J.~Beau,
``Exact Weil-Block Projective Capacity on Heisenberg Graphs: Resolving the Final Obstruction to $\delta$ Extraction,''
O12 spectral admissibility subprogram paper, Zenodo (2026).
DOI: \href{https://doi.org/10.5281/zenodo.19194798}{10.5281/zenodo.19194798}.

% 013
\bibitem{Beau2026a17}
J.~Beau,
``Asymptotic Stability of Exact Weil-Block Capacity on Heisenberg Graphs: Extended Prime Range, Variance Reduction,
and Requalification of the \(\delta \to \beta^*\) Tension,''
O13 spectral admissibility subprogram paper, Zenodo (2026).
DOI: \href{https://doi.org/10.5281/zenodo.19210831}{10.5281/zenodo.19210831}.

% 14
\bibitem{Beau2026a18}
J.~Beau,
``Observable-Class Mismatch and the Corrected \(\delta \to \beta^*\) Relation on Heisenberg Graphs:
Block Normalisation, Central-Phase Contribution,
and Structural Correspondence,''
O14 spectral admissibility subprogram paper, Zenodo (2026).
DOI: \href{https://doi.org/10.5281/zenodo.19226272}{10.5281/zenodo.19226272}.


    \bibitem{Wu1957}
    C.~S.~Wu, E.~Ambler, R.~W.~Hayward, D.~D.~Hoppes, R.~P.~Hudson,
    ``Experimental Test of Parity Conservation in Beta Decay,''
    Phys.\ Rev.\ \textbf{105}, 1413 (1957).

    \bibitem{PDG}
    S.~Navas \emph{et al.} (Particle Data Group),
    ``Review of Particle Physics,''
    Phys.\ Rev.\ D \textbf{110}, 030001 (2024).

    \bibitem{Patrizii2015}
    L.~Patrizii, M.~Spurio,
    ``Status of Searches for Magnetic Monopoles,''
    Ann.\ Rev.\ Nucl.\ Part.\ Sci.\ \textbf{65}, 279 (2015).

    \bibitem{Dirac1931}
    P.~A.~M.~Dirac,
    ``Quantised Singularities in the Electromagnetic Field,''
    Proc.\ Roy.\ Soc.\ London A \textbf{133}, 60 (1931).

    \bibitem{tHooft1974}
    G.~'t~Hooft,
    ``Magnetic Monopoles in Unified Gauge Theories,''
    Nucl.\ Phys.\ B \textbf{79}, 276 (1974).

    \bibitem{Polyakov1974}
    A.~M.~Polyakov,
    ``Particle Spectrum in Quantum Field Theory,''
    JETP Lett.\ \textbf{20}, 194 (1974).

  \end{thebibliography}

\end{document}
